\section{Related Work}\label{sec:related}

\emph{Labels.} Label and flow-construction errors in NIDS benchmarks are well documented for CIC-IDS2017 and CSE-CIC-IDS2018~\cite{engelen2021troubleshooting,liu2022error}, and correcting them changes detection scores~\cite{lanvin2023errors}. Flood et al.\ catalogue design smells of benchmark datasets~\cite{flood2024smells}, and models trained on such benchmarks can rely on shortcuts rather than attack behaviour~\cite{jacobs2022emperor}. Of 36 network security datasets analysed by Garcia and Valeros, 14 assign labels in file or directory names, the most common single practice~\cite{garcia2023labeling}. For CIC-IoT2023~\cite{neto2023ciciot2023}, Ramkumar et al.\ observed that the captures mix benign and attack traffic and re-filtered them per device by MAC address, but did not quantify the mixing~\cite{ramkumar2024diagnosing}; Vitorino et al.\ re-extract CIC-IoT2023 flows and assume that each capture contains only its named attack~\cite{vitorino2025revisiting}. Kostas et al.\ note that the attack pcaps of IoT-NID also contain benign packets, and that the derived IoTID20 flows label all of them as attacks~\cite[App.~A.3]{kostas2024iotgem}. Background traffic in attack captures is therefore not new; we add its share per capture and the cost of file-name labels on captures the model never saw.

\emph{Splits.} Random splits of time-ordered data are a known pitfall in security machine learning~\cite{arp2022dos,pendlebury2019tesseract,apruzzese2023sok}, time-aware evaluation has been applied to NIDS~\cite{andresini2021insomnia}, and tree classifiers degrade over time on MAWI flows~\cite{schraven2026mawiflow}. On IoT-23, Lyon et al.\ train on earlier months and find performance degrading as new malware types appear~\cite{lyon2026iot23drift}. In a time-shift experiment, Traffic-MoE uses a per-class chronological split (first 40\% of each class's time span for training, last 40\% for testing) on CIC-IoT2023 and CICIoMT2024, using byte-level flow tokens~\cite{zhou2026trafficmoe}; BRIDGE uses an early/late split in a pooled benchmark~\cite{bridge2026tchnet}; and Rasheed and Alobaedy find that drift causes no consistent decrease in performance but ``specific behaviour changes, such as total blindness, sensitivity inflation and relatively stable detection''~\cite{rasheed2026drift}. IoTGeM trains on telnet and tests on RTSP brute force, by design, and finds that the variants do not generalise~\cite{kostas2024iotgem}. Our split is within each capture, and we trace the gap per class and host pair.

\emph{Live detection and exporters.} Learned detectors on IoT gateways are common~\cite{realguard2022}, as are Mininet/Open vSwitch~\cite{lantz2010mininet,pfaff2015ovs} prototypes that block flagged hosts. Export timeouts delay flow-based detection~\cite{hofstede2013realtime}, the choice of expiration timeouts changes the flow records and can severely affect detection~\cite{velan2020impact}, and the exporter changes detection results~\cite{pinto2025flowexporter,pinto2025hera} and what a model sees~\cite{vitorino2025revisiting}. Models trained on complete flows lose up to 30\% precision and recall, under certain conditions, when tested on partial flows~\cite{pekar2025early}, and a streaming replay of graph-based NIDS found that none of 18 detector--dataset pairs could alert as events arrive~\cite{zhao2026gidseval}. We add a check that a live gateway classifies the flows of the offline evaluation.

\emph{Response.} Blocking via OpenFlow/OVS drop rules is standard in SDN DDoS mitigation~\cite{bawany2017ddos,yan2016sdn}, and a recent SoK reviews how machine-learning SDN defences are evaluated~\cite{orsdemir2026sok}. Few works we found report false-block counts on replayed traffic.
